% BHU PYQ -- Manually transcribed & error-corrected
% Paper: BCH-311: Advanced Company Accounts
% Exam: B.Com. (Hons.) Semester V Examination, 2013-14

\documentclass[11pt,a4paper]{article}
\usepackage[a4paper,top=2.5cm,bottom=2.5cm,left=2.8cm,right=2.8cm,headheight=14pt]{geometry}
\usepackage{amsmath,amssymb}
\usepackage[T1]{fontenc}
\usepackage[utf8]{inputenc}
\usepackage{lmodern,microtype,fancyhdr,enumitem,hyperref,lastpage,parskip}

\hypersetup{
    colorlinks=true,
    urlcolor=black,
    linkcolor=black,
    pdftitle={BCH-311: Advanced Company Accounts -- BHU B.Com. (Hons.) Sem V 2013-14}
}

\pagestyle{fancy}
\fancyhf{}
\renewcommand{\headrulewidth}{0.4pt}
\renewcommand{\footrulewidth}{0.4pt}
\fancyhead[L]{\small\textit{Banaras Hindu University}}
\fancyhead[R]{\small\textit{BCH-311: Advanced Company Accounts}}
\fancyfoot[C]{\small Page \thepage\ of \pageref{LastPage}}

\newlist{parts}{enumerate}{2}
\setlist[parts,1]{label=(\roman*),leftmargin=*,itemsep=5pt,topsep=3pt}
\setlist[parts,2]{label=(\alph*),leftmargin=*,itemsep=3pt,topsep=2pt}
\newcommand{\pts}[1]{\hfill\textit{[#1]}}

\begin{document}

\begin{center}
  {\large\bfseries BANARAS HINDU UNIVERSITY}\\[2pt]
  {\normalsize B.Com. (Hons.) --- Semester V Examination, 2013-14}\\[6pt]
  \rule{0.8\linewidth}{0.5pt}\\[6pt]
  {\large\bfseries Paper No. BCH-311: Advanced Company Accounts}\\[4pt]
  {\normalsize Time: 3 Hours\hfill Maximum Marks: 70}
\end{center}

\rule{\linewidth}{0.4pt}

\smallskip
\noindent\textit{Instructions: Answer all questions. All questions carry equal marks (14 marks each).}

\bigskip

\noindent\textbf{Question 1.} \\
The summarized Balance Sheets as on 31st March, 2013 of ABC Ltd. and XYZ Ltd. are as follows:\pts{14}
\begin{center}
\small
\renewcommand{\arraystretch}{1.2}
\begin{tabular}{|l|r|r|l|r|r|}
\hline
\textbf{Liabilities} & \textbf{ABC Ltd. (Rs.)} & \textbf{XYZ Ltd. (Rs.)} & \textbf{Assets} & \textbf{ABC Ltd. (Rs.)} & \textbf{XYZ Ltd. (Rs.)} \\ \hline
Share Capital (Rs. 100 each) & 15,00,000 & 5,00,000 & Goodwill & -- & 1,00,000 \\
Capital Reserve & -- & 50,000 & Building & 6,00,000 & -- \\
General Reserve & 2,00,000 & 60,000 & Plant \& Machinery & 6,50,000 & 4,20,000 \\
Profit \& Loss A/c & 1,20,000 & -- & Furniture & 10,000 & 5,000 \\
Sundry Creditors & 2,40,000 & 3,95,000 & Stock & 3,80,000 & 2,10,000 \\
& & & Debtors & 2,30,000 & 1,80,000 \\
& & & Cash \& Bank & 1,90,000 & 15,000 \\
& & & New Project Exp. & -- & 75,000 \\ \hline
\textbf{Total} & \textbf{20,60,000} & \textbf{10,05,000} & \textbf{Total} & \textbf{20,60,000} & \textbf{10,05,000} \\ \hline
\end{tabular}
\end{center}

XYZ Ltd. was amalgamated by ABC Ltd. on 1st April, 2013 on the following terms:
\begin{enumerate}[label=(\roman*),itemsep=2pt,topsep=2pt]
  \item Fixed assets other than goodwill were valued at Rs. 5,00,000 including Rs. 6,000 for furniture.
  \item Stock to be reduced by Rs. 20,000 and debtors by 5\%.
  \item ABC Ltd. assumed the liabilities of XYZ Ltd.
  \item The new projects to be valued at Rs. 95,000.
  \item The shareholders of XYZ Ltd. to receive cash payment of Rs. 30 per share plus 4 equity shares of ABC Ltd. for every five shares held.
  \item XYZ Ltd. to pay the liquidation expenses of ABC Ltd. (or XYZ Ltd.) amounting to Rs. 6,000.
\end{enumerate}
Pass the necessary journal entries in the books of XYZ Ltd. and prepare the Balance Sheet of ABC Ltd. after amalgamation.

\centerline{\textbf{OR}}

What is amalgamation of Companies? Discuss the important journal entries to be passed in the books of transferor and transferee.\pts{14}

\medskip\rule{\linewidth}{0.2pt}

\pagebreak

\noindent\textbf{Question 2.} \\
The following is the Balance Sheet of PQR Ltd. as on 31st March, 2013:\pts{14}
\begin{center}
\small
\renewcommand{\arraystretch}{1.2}
\begin{tabular}{|l|r|l|r|}
\hline
\textbf{Liabilities} & \textbf{Amount (Rs.)} & \textbf{Assets} & \textbf{Amount (Rs.)} \\ \hline
1,00,000 Equity shares of Rs. 10 each & 10,00,000 & Land & 1,00,000 \\
Sundry Creditors & 1,73,000 & Plant and Machinery & 2,30,000 \\
& & Furniture \& Fittings & 68,000 \\
& & Stock & 1,50,000 \\
& & Debtors & 70,000 \\
& & Cash at Bank & 5,000 \\
& & Profit \& Loss A/c & 5,50,000 \\ \hline
\textbf{Total} & \textbf{11,73,000} & \textbf{Total} & \textbf{11,73,000} \\ \hline
\end{tabular}
\end{center}

The terms of the approved scheme of Reconstruction are as follows:
\begin{enumerate}[label=(\roman*),itemsep=2pt,topsep=2pt]
  \item The equity shares to be reduced to Rs. 4 per share.
  \item Plant and machinery to be written down to Rs. 1,50,000.
  \item Stock to be revalued at Rs. 1,40,000.
  \item Provision for Bad debts to be made for Rs. 2,000.
  \item Land to be revalued at Rs. 1,42,000.
\end{enumerate}
Show the necessary journal entries for internal reconstruction and prepare the reconstructed Balance Sheet.

\centerline{\textbf{OR}}

Define and explain the following terms in the context of holding companies:\pts{14}
\begin{parts}
  \item Subsidiary Company
  \item Holding Company
  \item Minority interest
  \item Consolidated Balance Sheet
\end{parts}

\medskip\rule{\linewidth}{0.2pt}

\pagebreak

\noindent\textbf{Question 3.} \\
Prepare a Consolidated Balance Sheet of X Ltd. and its subsidiary Y Ltd. from the following Balance Sheets as on 31st March, 2013:\pts{14}
\begin{center}
\small
\renewcommand{\arraystretch}{1.2}
\begin{tabular}{|l|r|r|l|r|r|}
\hline
\textbf{Liabilities} & \textbf{X Ltd. (Rs.)} & \textbf{Y Ltd. (Rs.)} & \textbf{Assets} & \textbf{X Ltd. (Rs.)} & \textbf{Y Ltd. (Rs.)} \\ \hline
Share Capital (Rs. 10 each) & 25,00,000 & 6,00,000 & Land \& Building & 6,40,000 & 2,00,000 \\
General Reserve & 3,50,000 & 1,20,000 & Plant \& Machinery & 12,60,000 & 3,40,000 \\
Profit \& Loss A/c & 2,50,000 & 1,80,000 & Furniture \& Fittings & 1,40,000 & 60,000 \\
Sundry Creditors & 3,50,000 & 1,00,000 & Investment in Y Ltd. & & \\
& & & (40,000 shares) & 5,00,000 & -- \\
& & & Stock & 4,10,000 & 2,50,000 \\
& & & Debtors & 3,80,000 & 1,00,000 \\
& & & Cash and Bank & 1,20,000 & 50,000 \\ \hline
\textbf{Total} & \textbf{34,50,000} & \textbf{10,00,000} & \textbf{Total} & \textbf{34,50,000} & \textbf{10,00,000} \\ \hline
\end{tabular}
\end{center}
On the date of acquisition by X Ltd. of its holding of 40,000 shares in Y Ltd., Y Ltd. had undistributed profit and reserves amounting to Rs. 1,00,000.

\centerline{\textbf{OR}}

What do you mean by liquidation of companies? Prepare a liquidator's final statement of Account with imaginary figures.\pts{14}

\medskip\rule{\linewidth}{0.2pt}

\pagebreak

\noindent\textbf{Question 4.} \\
The following is the Balance Sheet of RST Ltd. as on 1st April, 2013:\pts{14}
\begin{center}
\small
\renewcommand{\arraystretch}{1.2}
\begin{tabular}{|l|r|l|r|}
\hline
\textbf{Liabilities} & \textbf{Amount (Rs.)} & \textbf{Assets} & \textbf{Amount (Rs.)} \\ \hline
24,000 Equity shares of Rs. 10 each & 2,40,000 & Machinery & 90,000 \\
Debentures & 1,50,000 & Leasehold Property & 1,20,000 \\
Bank overdraft & 54,000 & Stock & 3,000 \\
Creditors & 60,000 & Debtors & 1,50,000 \\
& & Investment & 18,000 \\
& & Cash & 3,000 \\
& & Profit \& Loss A/c & 1,20,000 \\ \hline
\textbf{Total} & \textbf{5,04,000} & \textbf{Total} & \textbf{5,04,000} \\ \hline
\end{tabular}
\end{center}

RST Ltd. went into voluntary liquidation on the Balance Sheet date. The assets were realized as follows:
\begin{center}
\begin{tabular}{lr}
  Machinery & Rs. 1,80,000 \\
  Leasehold Property & Rs. 2,18,000 \\
  Investment & Rs. 12,000 \\
  Stock & Rs. 6,000 \\
  Debtors & Rs. 1,40,000 \\
\end{tabular}
\end{center}
The bank overdraft is secured by the deposit of the documents of leasehold property. Preferential creditors amounted to Rs. 3,000, which were not included in creditors of Rs. 60,000. Prepare the Statement of Affairs or Liquidator's Final Statement of Account.

\centerline{\textbf{OR}}

What is meant by Bonus Share? State the accounting treatment of Bonus Shares issue.\pts{14}

\medskip\rule{\linewidth}{0.2pt}

\noindent\textbf{Question 5.} \\
Write short notes on the following:\pts{7+7}
\begin{parts}
  \item Employee stock option Scheme
  \item Buy-Back of Shares
\end{parts}

\centerline{\textbf{OR}}

What is valuation of goodwill? Explain the various methods of valuation of goodwill of a company.\pts{14}

\vfill
\rule{\linewidth}{0.4pt}
\begin{center}\small
\href{https://bhu-exam.vercel.app/}{Click here to visit our website}\\[4pt]
\href{https://me-aryan.vercel.app/}{Click here to visit our contributor}\\[4pt]
\href{https://quashan-me.vercel.app/}{Click here to visit our contributor}
\end{center}

\end{document}
